% Options for packages loaded elsewhere
\PassOptionsToPackage{unicode}{hyperref}
\PassOptionsToPackage{hyphens}{url}
\documentclass[
]{article}
\usepackage{xcolor}
\usepackage[margin=1in]{geometry}
\usepackage{amsmath,amssymb}
\setcounter{secnumdepth}{-\maxdimen} % remove section numbering
\usepackage{iftex}
\ifPDFTeX
  \usepackage[T1]{fontenc}
  \usepackage[utf8]{inputenc}
  \usepackage{textcomp} % provide euro and other symbols
\else % if luatex or xetex
  \usepackage{unicode-math} % this also loads fontspec
  \defaultfontfeatures{Scale=MatchLowercase}
  \defaultfontfeatures[\rmfamily]{Ligatures=TeX,Scale=1}
\fi
\usepackage{lmodern}
\ifPDFTeX\else
  % xetex/luatex font selection
\fi
% Use upquote if available, for straight quotes in verbatim environments
\IfFileExists{upquote.sty}{\usepackage{upquote}}{}
\IfFileExists{microtype.sty}{% use microtype if available
  \usepackage[]{microtype}
  \UseMicrotypeSet[protrusion]{basicmath} % disable protrusion for tt fonts
}{}
\makeatletter
\@ifundefined{KOMAClassName}{% if non-KOMA class
  \IfFileExists{parskip.sty}{%
    \usepackage{parskip}
  }{% else
    \setlength{\parindent}{0pt}
    \setlength{\parskip}{6pt plus 2pt minus 1pt}}
}{% if KOMA class
  \KOMAoptions{parskip=half}}
\makeatother
\usepackage{longtable,booktabs,array}
\newcounter{none} % for unnumbered tables
\usepackage{calc} % for calculating minipage widths
% Correct order of tables after \paragraph or \subparagraph
\usepackage{etoolbox}
\makeatletter
\patchcmd\longtable{\par}{\if@noskipsec\mbox{}\fi\par}{}{}
\makeatother
% Allow footnotes in longtable head/foot
\IfFileExists{footnotehyper.sty}{\usepackage{footnotehyper}}{\usepackage{footnote}}
\makesavenoteenv{longtable}
\usepackage{graphicx}
\makeatletter
\newsavebox\pandoc@box
\newcommand*\pandocbounded[1]{% scales image to fit in text height/width
  \sbox\pandoc@box{#1}%
  \Gscale@div\@tempa{\textheight}{\dimexpr\ht\pandoc@box+\dp\pandoc@box\relax}%
  \Gscale@div\@tempb{\linewidth}{\wd\pandoc@box}%
  \ifdim\@tempb\p@<\@tempa\p@\let\@tempa\@tempb\fi% select the smaller of both
  \ifdim\@tempa\p@<\p@\scalebox{\@tempa}{\usebox\pandoc@box}%
  \else\usebox{\pandoc@box}%
  \fi%
}
% Set default figure placement to htbp
\def\fps@figure{htbp}
\makeatother
\setlength{\emergencystretch}{3em} % prevent overfull lines
\providecommand{\tightlist}{%
  \setlength{\itemsep}{0pt}\setlength{\parskip}{0pt}}
\usepackage{bookmark}
\IfFileExists{xurl.sty}{\usepackage{xurl}}{} % add URL line breaks if available
\urlstyle{same}
\hypersetup{
  pdftitle={Preliminary Results},
  pdfauthor={Juliana Balluffi-Fry},
  hidelinks,
  pdfcreator={LaTeX via pandoc}}

\title{Preliminary Results}
\author{Juliana Balluffi-Fry}
\date{2026-09-18}

\begin{document}
\maketitle

Results

General Summary

Across the entire study, squirrels experienced an average local density
of 2.3 ± 1.2 squirrels/ha. Food supplementation increased the local
density around mothers on food add grids at a rate of 0.12 squirrel per
hectare per year over the 12 years of its food experiment. Over the
entire study, mothers on food supplemented grids experienced higher
local densities, with an average of 2.59 squirrels/ha, compared to those
on control grids, that had on average 2.15 squirrels/ha. Mothers on
rattle treatment grids had the lowest densities at 1.61 squirrels/ha
(Figure 1A). However, there was substantial natural variation in local
densities over time (Figure 2B). Pups grew at a rate of 1.8 ± 0.5 g/day,
with litters from bucket middens have slightly higher average growth
rates at 1.8 ± 0.4 g/day (Figure 1C). Litters had 3 pups on average, but
could range anywhere from 1 to 7 pups (Figure 1D). Based on data from
control grids, with every pup increase in litter size, litter growth
rate decreased 0.16 g/day (Figure 2). The most common age for mothers
was 2 years, and litters were born on average at 116 ± 24 Julian days
into the year.

Tables

\begin{longtable}[]{@{}llrrrl@{}}
\caption{Table 1. Sample size of juveniles, litters, and female red
squirrels and the years of those sample sizes, by grid and
grid-treatment.}\tabularnewline
\toprule\noalign{}
Grid & Grid\_treatment & N\_juveniles & N\_litters & N\_females &
Years \\
\midrule\noalign{}
\endfirsthead
\toprule\noalign{}
Grid & Grid\_treatment & N\_juveniles & N\_litters & N\_females &
Years \\
\midrule\noalign{}
\endhead
\bottomrule\noalign{}
\endlastfoot
KL & Control & 2214 & 720 & 415 & 1989-2025 \\
SU & Control & 1461 & 490 & 290 & 1989-2025 \\
LL & Food add & 506 & 142 & 90 & 2006-2011 \\
JO & Food add & 460 & 129 & 85 & 2007-2012 \\
AG & Food add & 1387 & 436 & 257 & 2005-2017 \\
AG & Rattle & 15 & 4 & 4 & 2010-2010 \\
JO & Rattle & 150 & 47 & 41 & 2018-2022 \\
BT & Control & 38 & 12 & 12 & 2010-2021 \\
LL & Rattle & 45 & 11 & 11 & 2010-2010 \\
BT & Rattle & 14 & 4 & 4 & 2010-2010 \\
\end{longtable}

\begin{longtable}[]{@{}llrl@{}}
\caption{Table 2. Control Model Results.}\tabularnewline
\toprule\noalign{}
Response & Effect & Estimate & CI \\
\midrule\noalign{}
\endfirsthead
\toprule\noalign{}
Response & Effect & Estimate & CI \\
\midrule\noalign{}
\endhead
\bottomrule\noalign{}
\endlastfoot
GR & Intercept & 0.08 & (-0.03, 0.19) \\
GR & density\_std & 0.00 & (-0.06, 0.06) \\
GR & cones\_std & 0.00 & (-0.06, 0.05) \\
GR & mother\_age\_std & -0.10 & (-0.15, -0.04) \\
GR & bjulian\_std & -0.14 & (-0.19, -0.08) \\
GR & mastyes & -0.03 & (-0.2, 0.14) \\
GR & sexratio & -0.15 & (-0.32, 0.02) \\
LS & Intercept & -0.06 & (-0.17, 0.05) \\
LS & density\_std & -0.04 & (-0.1, 0.02) \\
LS & cones\_std & 0.01 & (-0.05, 0.07) \\
LS & mother\_age\_std & 0.03 & (-0.03, 0.09) \\
LS & bjulian\_std & -0.05 & (-0.11, 0.01) \\
LS & mastyes & 0.00 & (-0.18, 0.17) \\
LS & sexratio & 0.10 & (-0.08, 0.28) \\
rescor & NA & -0.34 & (-0.39, -0.29) \\
sd(GR) & gridyear & 0.02 & (0, 0.07) \\
sd(GR) & mother\_byear & 0.05 & (0, 0.13) \\
sd(LS) & gridyear & 0.02 & (0, 0.07) \\
sd(LS) & mother\_byear & 0.07 & (0, 0.17) \\
\end{longtable}

\begin{longtable}[]{@{}llrl@{}}
\caption{Table 3. Food Model Results}\tabularnewline
\toprule\noalign{}
Response & Effect & Estimate & CI \\
\midrule\noalign{}
\endfirsthead
\toprule\noalign{}
Response & Effect & Estimate & CI \\
\midrule\noalign{}
\endhead
\bottomrule\noalign{}
\endlastfoot
GR & Intercept & 0.06 & (-0.09, 0.21) \\
GR & treatment\_indNoBucket & -0.12 & (-0.3, 0.05) \\
GR & treatment\_indBucket & 0.11 & (-0.03, 0.26) \\
GR & density\_std & -0.05 & (-0.15, 0.05) \\
GR & cones\_std & 0.03 & (-0.04, 0.1) \\
GR & mother\_age\_std & -0.02 & (-0.08, 0.04) \\
GR & bjulian\_std & -0.17 & (-0.23, -0.11) \\
GR & mastyes & 0.01 & (-0.15, 0.17) \\
GR & sexratio & -0.19 & (-0.39, -0.01) \\
GR & treatment\_indNoBucket:density\_std & 0.12 & (-0.03, 0.27) \\
GR & treatment\_indBucket:density\_std & 0.02 & (-0.13, 0.16) \\
LS & Intercept & -0.08 & (-0.23, 0.06) \\
LS & treatment\_indNoBucket & 0.06 & (-0.11, 0.24) \\
LS & treatment\_indBucket & 0.05 & (-0.09, 0.19) \\
LS & density\_std & -0.05 & (-0.15, 0.05) \\
LS & cones\_std & 0.03 & (-0.04, 0.1) \\
LS & mother\_age\_std & 0.07 & (0.01, 0.13) \\
LS & bjulian\_std & -0.10 & (-0.16, -0.04) \\
LS & mastyes & -0.03 & (-0.19, 0.13) \\
LS & sexratio & 0.16 & (-0.03, 0.35) \\
LS & treatment\_indNoBucket:density\_std & -0.06 & (-0.22, 0.09) \\
LS & treatment\_indBucket:density\_std & -0.08 & (-0.23, 0.06) \\
rescor & NA & -0.30 & (-0.36, -0.25) \\
sd(GR) & gridyear & 0.03 & (0, 0.07) \\
sd(GR) & mother\_byear & 0.10 & (0.01, 0.22) \\
sd(LS) & gridyear & 0.03 & (0, 0.08) \\
sd(LS) & mother\_byear & 0.07 & (0, 0.18) \\
\end{longtable}

\begin{longtable}[]{@{}llrl@{}}
\caption{Table 3. Cache Model Results}\tabularnewline
\toprule\noalign{}
Response & Effect & Estimate & CI \\
\midrule\noalign{}
\endfirsthead
\toprule\noalign{}
Response & Effect & Estimate & CI \\
\midrule\noalign{}
\endhead
\bottomrule\noalign{}
\endlastfoot
GR & Intercept & 0.14 & (-0.11, 0.39) \\
GR & cache\_catmedium & 0.02 & (-0.2, 0.25) \\
GR & cache\_catlarge & 0.00 & (-0.26, 0.27) \\
GR & density\_std & -0.18 & (-0.43, 0.07) \\
GR & cones\_std & 0.00 & (-0.09, 0.09) \\
GR & mother\_age\_std & -0.05 & (-0.16, 0.06) \\
GR & bjulian\_std & -0.16 & (-0.26, -0.06) \\
GR & mastyes & -0.03 & (-0.36, 0.29) \\
GR & sexratio & -0.28 & (-0.57, 0.01) \\
GR & cache\_catmedium:density\_std & 0.16 & (-0.12, 0.44) \\
GR & cache\_catlarge:density\_std & 0.22 & (-0.08, 0.53) \\
LS & Intercept & -0.16 & (-0.42, 0.1) \\
LS & cache\_catmedium & 0.00 & (-0.23, 0.23) \\
LS & cache\_catlarge & 0.01 & (-0.26, 0.28) \\
LS & density\_std & 0.13 & (-0.13, 0.38) \\
LS & cones\_std & -0.03 & (-0.12, 0.06) \\
LS & mother\_age\_std & -0.03 & (-0.14, 0.08) \\
LS & bjulian\_std & -0.09 & (-0.19, 0) \\
LS & mastyes & 0.02 & (-0.31, 0.35) \\
LS & sexratio & 0.23 & (-0.06, 0.53) \\
LS & cache\_catmedium:density\_std & -0.05 & (-0.33, 0.24) \\
LS & cache\_catlarge:density\_std & -0.15 & (-0.46, 0.15) \\
rescor & NA & -0.36 & (-0.44, -0.27) \\
sd(GR) & gridyear & 0.05 & (0, 0.15) \\
sd(GR) & mother\_byear & 0.08 & (0, 0.24) \\
sd(LS) & gridyear & 0.05 & (0, 0.13) \\
sd(LS) & mother\_byear & 0.15 & (0.01, 0.34) \\
\end{longtable}

\begin{longtable}[]{@{}llrl@{}}
\caption{Table 4. Rattle Model Results}\tabularnewline
\toprule\noalign{}
Response & Effect & Estimate & CI \\
\midrule\noalign{}
\endfirsthead
\toprule\noalign{}
Response & Effect & Estimate & CI \\
\midrule\noalign{}
\endhead
\bottomrule\noalign{}
\endlastfoot
GR & Intercept & 0.00 & (-0.35, 0.35) \\
GR & treatment\_grRattle & -0.09 & (-0.48, 0.3) \\
GR & density\_std & 0.02 & (-0.21, 0.26) \\
GR & cones\_std & -0.02 & (-0.19, 0.16) \\
GR & bjulian\_std & 0.04 & (-0.1, 0.17) \\
GR & mastyes & -0.09 & (-0.5, 0.31) \\
GR & sexratio & 0.09 & (-0.35, 0.52) \\
GR & treatment\_grRattle:density\_std & -0.07 & (-0.53, 0.38) \\
LS & Intercept & 0.14 & (-0.2, 0.48) \\
LS & treatment\_grRattle & -0.13 & (-0.52, 0.25) \\
LS & density\_std & -0.08 & (-0.33, 0.16) \\
LS & cones\_std & -0.02 & (-0.18, 0.15) \\
LS & bjulian\_std & -0.14 & (-0.27, 0) \\
LS & mastyes & -0.13 & (-0.53, 0.27) \\
LS & sexratio & -0.26 & (-0.69, 0.17) \\
LS & treatment\_grRattle:density\_std & -0.31 & (-0.76, 0.13) \\
rescor & NA & -0.31 & (-0.42, -0.18) \\
sd(GR) & gridyear & 0.07 & (0, 0.22) \\
sd(GR) & mother\_byear & 0.14 & (0.01, 0.4) \\
sd(LS) & gridyear & 0.07 & (0, 0.22) \\
sd(LS) & mother\_byear & 0.13 & (0.01, 0.37) \\
\end{longtable}

\begin{longtable}[]{@{}
  >{\raggedright\arraybackslash}p{(\linewidth - 6\tabcolsep) * \real{0.1324}}
  >{\raggedright\arraybackslash}p{(\linewidth - 6\tabcolsep) * \real{0.5147}}
  >{\raggedleft\arraybackslash}p{(\linewidth - 6\tabcolsep) * \real{0.1324}}
  >{\raggedright\arraybackslash}p{(\linewidth - 6\tabcolsep) * \real{0.2206}}@{}}
\caption{Table 5. Rattle Model Results with Chickadee
test.}\tabularnewline
\toprule\noalign{}
\begin{minipage}[b]{\linewidth}\raggedright
Response
\end{minipage} & \begin{minipage}[b]{\linewidth}\raggedright
Effect
\end{minipage} & \begin{minipage}[b]{\linewidth}\raggedleft
Estimate
\end{minipage} & \begin{minipage}[b]{\linewidth}\raggedright
CI
\end{minipage} \\
\midrule\noalign{}
\endfirsthead
\toprule\noalign{}
\begin{minipage}[b]{\linewidth}\raggedright
Response
\end{minipage} & \begin{minipage}[b]{\linewidth}\raggedright
Effect
\end{minipage} & \begin{minipage}[b]{\linewidth}\raggedleft
Estimate
\end{minipage} & \begin{minipage}[b]{\linewidth}\raggedright
CI
\end{minipage} \\
\midrule\noalign{}
\endhead
\bottomrule\noalign{}
\endlastfoot
GR & Intercept & 0.00 & (-0.34, 0.36) \\
GR & treatment\_indChickadee & -0.09 & (-0.94, 0.74) \\
GR & treatment\_indRattle & -0.08 & (-0.48, 0.31) \\
GR & density\_std & 0.05 & (-0.2, 0.3) \\
GR & cones\_std & -0.02 & (-0.2, 0.16) \\
GR & bjulian\_std & 0.04 & (-0.1, 0.18) \\
GR & mastyes & -0.11 & (-0.54, 0.31) \\
GR & sexratio & 0.08 & (-0.35, 0.52) \\
GR & treatment\_indChickadee:density\_std & -0.27 & (-1.14, 0.58) \\
GR & treatment\_indRattle:density\_std & -0.10 & (-0.56, 0.36) \\
LS & Intercept & 0.14 & (-0.2, 0.49) \\
LS & treatment\_indChickadee & -0.09 & (-0.92, 0.74) \\
LS & treatment\_indRattle & -0.14 & (-0.54, 0.24) \\
LS & density\_std & -0.08 & (-0.33, 0.18) \\
LS & cones\_std & -0.01 & (-0.18, 0.16) \\
LS & bjulian\_std & -0.14 & (-0.27, 0) \\
LS & mastyes & -0.11 & (-0.54, 0.3) \\
LS & sexratio & -0.26 & (-0.68, 0.17) \\
LS & treatment\_indChickadee:density\_std & -0.07 & (-0.89, 0.77) \\
LS & treatment\_indRattle:density\_std & -0.31 & (-0.78, 0.13) \\
rescor & NA & -0.31 & (-0.42, -0.18) \\
sd(GR) & gridyear & 0.08 & (0, 0.22) \\
sd(GR) & mother\_byear & 0.14 & (0.01, 0.4) \\
sd(LS) & gridyear & 0.08 & (0, 0.23) \\
sd(LS) & mother\_byear & 0.14 & (0.01, 0.38) \\
\end{longtable}

Figure Captions

Figure 1. Summary of the annual changes in A) squirrel densities
(squirrels/ha) and B) spruce cone production (log cones/tree) over the
37-long study in Kluane, Yukon. Lines representing squirrel densities
(panel A) are colored by treatment of the study area (grey = food add,
red = food add, fuchsia = rattle playbacks). Dashed vertical lines of
the same colors delineate the time frames of each experiment, the food
supplementation experiment (red) and the rattle playback experiment
(fuchsia). Diamond points along the line representing cone production
(panel B) indicates that the associated year was categorized as a mast
year. Lines represent annual means and shaded areas represent the
standard deviation.

Figure 2. Raw data from control grids (n = 1198) showing the overall
negative trade-off between litter size and litter growth rate (g/day).
Linear relationship is shown.

Figure 3. Best linear unbiased predictors from all three Bayesian
models, food (n = 1178), cache (n = 454), and rattle (n = 238), that
test how growth rate and litter size respond to local density. After all
fixed effects and random effects are accounted for, this is how litter
size and growth rate relate to one another in each scenario. Data points
represent litters and are colored based on either A) food treatment, B)
cache size, or C) rattle treatment. Linear relationships and 95\%
confidence levels are represented by lines and ellipses, respectively.

Figure 4. Predictions of the covariance (left column) and correlation
(right column) between litter growth rate (GR) and litter size (LS) in
response to local density based on three Bayesian models: Food model
(top row), Cache model (middle row), and Rattle model (bottom row). Each
model tested a different scenario or experiment. The Food model (n =
1178) tested how the interaction between a mother's local density and
her food treatment (control = grey, food add grid with a bucket = red,
food add grid with no bucket = blue) affects the relationship between
litter growth rate and litter size. The Cache model (n = 454) tested how
the interaction between a mother's local density and her cache size in
the previous autumn (small = yellow, medium = light green, large = dark
green) affects the relationship between litter growth rate and litter
size. The Rattle model (n = 238) tested how the interaction between a
mother's local density and her perceived density treatment (control =
grey, rattle playback = violet) affects the relationship between litter
growth rate and litter size.

Figure 5. Linear predictions for litter growth rate (top row) and litter
size (bottom row), separately, in response to local density based on
four Bayesian models: A) Control model, B) Food model, C) Cache model,
and D) Rattle model. Each model tested a different scenario or
experiment, with the control model (A) being the base model. B)
Predictions from the Food model (n = 1178) tested how the interaction
between a mother's local density and her food treatment (control = grey,
food add grid with a bucket = red, food add grid with no bucket = blue)
affects the relationship between litter growth rate and litter size. C)
Predictions from the Cache model (n = 454) that tested how the
interaction between a mother's local density and her cache size in the
previous autumn (small = yellow, medium = light green, large = dark
green) affects the relationship between litter growth rate and litter
size. D) Predictions from the Rattle model (n = 238) that tested how the
interaction between a mother's local density and her perceived density
treatment (control = grey, rattle playback = violet) affects the
relationship between litter growth rate and litter size.

Figures

\begin{figure}
\centering
\pandocbounded{\includegraphics[keepaspectratio,alt={Figure 1}]{Results_files/Output/Figures/annual_summary.jpeg}}
\caption{Figure 1}
\end{figure}

\begin{figure}
\centering
\pandocbounded{\includegraphics[keepaspectratio,alt={Figure 2}]{Results_files/Output/Figures/GR_littersize.jpeg}}
\caption{Figure 2}
\end{figure}

\begin{figure}
\centering
\pandocbounded{\includegraphics[keepaspectratio,alt={Figure 3}]{Results_files/Output/Figures/ellipse_predictions.jpeg}}
\caption{Figure 3}
\end{figure}

\begin{figure}
\centering
\pandocbounded{\includegraphics[keepaspectratio,alt={Figure 4}]{Results_files/Output/Figures/covariance_correlation_predictions.jpeg}}
\caption{Figure 4}
\end{figure}

\begin{figure}
\centering
\pandocbounded{\includegraphics[keepaspectratio,alt={Figure 5.}]{Results_files/Output/Figures/Bayesian_separate_predictions.jpeg}}
\caption{Figure 5.}
\end{figure}

\end{document}
